\begin{figure*}[t]
    \centering
    \begingroup
    \definecolor{responsehealthyframe}{HTML}{3C7D73}
    \definecolor{responsehealthyfill}{HTML}{F1F8F6}
    \definecolor{responsecollapseframe}{HTML}{B85245}
    \definecolor{responsecollapsefill}{HTML}{FCF3F1}
    \setlength{\fboxsep}{5pt}
    \setlength{\fboxrule}{0.55pt}
    \newcommand{\responseexample}[5]{%
        \begin{minipage}[t]{0.485\textwidth}
            \fcolorbox{#1}{#2}{%
                \parbox[t][25mm][t]{\dimexpr\linewidth-2\fboxsep-2\fboxrule\relax}{%
                    \raggedright
                    \scriptsize
                    \textbf{#3}\par
                    {\fontsize{6.8}{7.5}\selectfont\textbf{#4}}\par
                    \vspace{1.5pt}
                    #5
                }%
            }
        \end{minipage}%
    }

    \responseexample{responsehealthyframe}{responsehealthyfill}
        {(a) Teacher: coherent solution}{correct; 842 tokens}{%
        ``Number of \textbf{grand prize} outcomes: 1. \ldots\ Number of
        \textbf{prize} outcomes ($\geq 2$ matches): 115. \ldots
        $P(\text{Grand Prize}\mid\text{Prize})=1/115$.
        Thus, $m+n=1+115=116$.''}
    \hfill
    \responseexample{responsehealthyframe}{responsehealthyfill}
        {(b) $S_{\mathrm{INIT}}$: coherent solution}{correct; 771 tokens}{%
        ``Jen has exactly 2 of the chosen 4 numbers:
        $C(4,2)C(6,2)=90$. \ldots\ Exactly 3: 24; all 4: 1.
        So, the total number of ways to win a prize is
        $90+24+1=115$. \ldots
        $P(\text{grand prize}\mid\text{prize})=1/115$.
        \ldots\ \fbox{116}.''}

    \vspace{4pt}

    \responseexample{responsecollapseframe}{responsecollapsefill}
        {(c) $S_{\mathrm{INIT}}$, teacher-preferred: repeated reconsideration}{incorrect; 7,311 tokens}{%
        ``$17{,}499{,}840/661{,}680=26.5$.
        \textbf{Wait, this is incorrect.} The calculation is not correct.
        Instead, we should compute the number of valid sequences. Let's go
        back. \ldots\ But this is not correct. \textbf{Wait, for the last word}
        to be GY, \ldots\ \textbf{Wait, this is not correct.}''}
    \hfill
    \responseexample{responsecollapseframe}{responsecollapsefill}
        {(d) $S_{\mathrm{INIT}}$, teacher-preferred: mechanical repetition}{correct; 8,192 tokens; truncated}{%
        ``\textbf{Final Answer: \fbox{116}} \ldots\
        \textbf{Answer: \fbox{116}} \ldots\
        \textbf{Final Answer: \fbox{116}} \ldots\
        \textbf{Answer: \fbox{116}} \ldots''\par
        \vspace{2pt}
        The same answer line repeats 365 times until the generation length
        limit.}
    \endgroup
    \caption{
        \textbf{Illustrative response excerpts from the Qwen3-4B setting.}
        The teacher and a typical $S_{\mathrm{INIT}}$ response give coherent
        solutions, whereas the $S_{\mathrm{INIT}}$ responses the teacher most
        prefers repeatedly restart their reasoning or mechanically repeat an
        answer until the generation length limit. All four excerpts predate any
        OPD update. Ellipses mark omitted text. These examples illustrate
        behaviors and do not estimate their prevalence.
    }
    \label{fig:qualitative_response_examples}
\end{figure*}
